% Part: first-order-logic
% Chapter: beyond
% Section: intuitionistic-logic

\documentclass[../../../include/open-logic-section]{subfiles}

\begin{document}

\olfileid{fol}{byd}{il}

\olsection{Logika Intuisionistik}

Beda karo logika orde kapindho lan orde luwih dhuwur,
logika orde kapisan intuisionistik iku wujud
logika klasik kang diwatesi, kanthi tujuan
nggambarake cara nalar kang luwih "konstruktif".
Tuladha-tuladha ing ngisor iki bisa njlentrehake
sawetara panyebab kang ndhasari gagasan kasebut.

Upamakna ing sawijining dina ana wong kandha
marang sampeyan yen dheweke wis nemtokake
wilangan asli~$x$ kang nduweni sipat iki:
yen $x$ prima, hipotesis Riemann bener;
yen $x$ komposit, hipotesis Riemann luput.
Iki kabar apik!{} Apa hipotesis Riemann
bener utawa ora isih dadi salah siji
pitakonan gedhe kang durung kaudhari
ing matematika. Wong iku kaya-kaya wis
ngowahi masalah mau dadi pitakonan
petungan: apa wilangan tartamtu iku
prima utawa dudu.

Pira nilai ajaib saka $x$ iku?
Miturut katrangane, $x$ iku wilangan
asli kang padha karo $7$ yen
hipotesis Riemann bener, lan
padha karo $9$ yen ora bener.

Sampeyan nesu lan njaluk dhuwit
sampeyan dibalekake. Saka sudut
pandang klasik, katrangan ing
ndhuwur pancen nemtokake siji-sijine
nilai~$x$; nanging kang sejatine
sampeyan karepake yaiku nilai~$x$
kang diwenehake kanthi \emph{gamblang}.

Minangka tuladha liyane kang mbokmenawa
ora banget digawe-gawe, gatekna pitakonan
iki. Kita ngerti yen wilangan irasional
bisa dipangkatake nganggo pangkat rasional
lan ngasilake wilangan rasional.
Upamane, $\sqrt{2}^2 = 2$.
Kang durung cetha yaiku apa wilangan
irasional bisa dipangkatake nganggo
pangkat \emph{irasional} lan ngasilake
wilangan rasional. Teorema ing ngisor
iki mangsuli ya:

\begin{thm}
Ana wilangan irasional $a$ lan $b$
saéngga $a^b$ rasional.
\end{thm}

\begin{proof}
Gatekna $\sqrt{2}^{\sqrt{2}}$.
Yen iki rasional, prakara rampung:
kita bisa njupuk $a = b = \sqrt{2}$.
Yen ora, wilangan iki irasional.
Mula kita duwe
\[
(\sqrt{2}^{\sqrt{2}})^{\sqrt{2}} = \sqrt{2}^{\sqrt{2} \cdot
  \sqrt{2}} = \sqrt{2}^2 = 2,
\]
kang mesthi rasional. Dadi, ing kahanan
iki, jupuken $a$ minangka
$\sqrt{2}^{\sqrt{2}}$, lan
$b$ minangka~$\sqrt 2$.
\end{proof}

Apa iki dadi bukti kang sah?
Umume ahli matematika rumangsa
yen mangkono. Nanging ana bab
kang isih kurang marem: kita
wis mbuktekake anane pasangan
wilangan real kanthi sipat
tartamtu, tanpa bisa nyebutake
\emph{pasangan endi} kang dimaksud.
Asil kang padha bisa dibuktekake
kanthi cara kang menehi pasangan
$a$, $b$ \emph{ing bukti iku dhewe}:
jupuken $a = \sqrt{3}$ lan
$b = \log_3 4$. Banjur
\[
a^b = \sqrt{3}^{\log_3 4} = 3^{1/2 \cdot \log_3 4} = (3^{\log_3
  4})^{1/2} = 4^{1/2}= 2,
\]
amarga $3^{\log_3 x} = x$. Cathetan koreksi sumber OLPL-753: b uga kudu dibuktekake irasional. Yen logaritma basis telu saka papat padha pecahan p per q kanggo wilangan bulat positif p lan q, mula telu pangkat p padha papat pangkat q, kang mokal amarga faktor primane beda.

Logika intuisionistik dirancang
kanggo nggambarake cara nalar
kang ora ngidini langkah kaya
ing bukti kapisan. Kanggo
mbuktekake anane $x$ kang
nyukupi~$!A(x)$, kita kudu
menehi~$x$ tartamtu lan bukti
yen nilai iku nyukupi $!A$,
kaya ing bukti kapindho.
Kanggo mbuktekake yen $!A$
utawa $!B$ bener, kita kudu
bisa mbuktekake salah sijine.

Sacara formal, logika orde kapisan
intuisionistik dioleh kanthi
matesi sistem !!a{derivation}
kanggo logika orde kapisan
kanthi cara tartamtu.
Semono uga, ana versi
intuisionistik kanggo logika
orde kapindho utawa orde
luwih dhuwur. Saka sudut
pandang matematika, sistem
kasebut mung sistem deduktif
formal, nanging, kaya kang
wis kacathet, tujuane
nggambarake cara nalar
matematika tartamtu.
Cara nalar iki bisa dianggep
minangka cara kang dibenerake
dening pandangan filsafat
matematika tartamtu (kayata
intuisionisme Brouwer);
uga bisa dianggep minangka
cara nalar matematika kang
luwih "konkret" lan marem
(kaya konstruktivisme Bishop).
Wong uga bisa padha beda
panemu apa gambaran formal
iku cocog karo alesan informal.
Nanging apa wae pandangan
filsafat kang dianut,
logika intuisionistik bisa
ditliti minangka logika
kang diwenehake sacara
formal; lan amarga warna-warna
alesan, akeh ahli logika
matematika kang seneng
nliti logika iki.

Ana tafsiran konstruktif
informal kanggo panyambung
intuisionistik, kang lumrahe
diarani tafsiran BHK
(saka jeneng Brouwer,
Heyting, lan Kolmogorov).
Tafsiran iku kaya mangkene:
bukti $!A \land !B$
dumadi saka bukti $!A$
kang dipasangake karo
bukti $!B$; bukti
$!A \lor !B$ dumadi
saka bukti $!A$ utawa
bukti $!B$, bebarengan
karo katrangan gamblang
endi kang kelakon;
bukti $!A \lif !B$
iku prosedur kang ngowahi
bukti $!A$ dadi bukti~$!B$;
bukti $\lforall[x][!A(x)]$
iku prosedur kang ngasilake
bukti $!A(x)$ kanggo
saben nilai~$x$; lan bukti
$\lexists[x][!A(x)]$
dumadi saka nilai~$x$
bebarengan karo bukti
yen nilai iku nyukupi~$!A$.
Tafsiran iki uga bisa
diandharake kanthi istilah
komputasi kang diarani
"isomorfisme Curry--Howard"
utawa "paradigma
!!{formula}s-minangka-tipe":
!!a{formula} dianggep
nemtokake jinis tipe data,
dene bukti dadi objek
komputasi saka tipe data
kasebut kang nuduhake
yen !!{formula} kang
cocog iku bener.

Logika intuisionistik asring
dianggep logika klasik
"tanpa" hukum tengah
kang disingkirake.
Teorema ing ngisor iki
njlentrehake teges
pernyataan kasebut.
\begin{thm}
Sacara intuisionistik,
skema aksioma ing ngisor
iki padha kuwat:
\begin{enumerate}
\item $(\lnot !A \lif \lfalse) \lif !A$.
\item $!A \lor \lnot !A$
\item $\lnot \lnot !A \lif !A$
\end{enumerate}
\end{thm}
Nuduhake yen instansi saka
salah siji skema bisa dioleh
saka salah siji skema liyane
dadi gladhen kang becik
ing logika intuisionistik.

Sistem deduktif kapisan
kanggo logika proposisional
intuisionistik, kang diajokake
kanggo formalisasi intuisionisme
Brouwer, digawe kanthi mandiri
dening Kolmogorov, Glivenko,
lan Heyting. Formalisasi
kapisan logika orde kapisan
intuisionistik (lan saperangan
matematika intuisionistik)
asal saka Heyting. Sanadyan
sawetara skema sah ing logika
klasik nanging ora sah sacara
intuisionistik, akeh skema
liyane kang tetep sah.

\emph{Terjemahan negasi pindho}
njlentrehake sesambungan
wigati antarane logika klasik
lan intuisionistik.
Terjemahan iki ditemtokake
kanthi induksi kaya mangkene
(anggepen $!A^N$ minangka
terjemahan "intuisionistik"
saka !!{formula} klasik~$!A$):
\begin{align*}
!A^N & \ident \lnot\lnot !A \quad \text{kanggo !!{formula}s atomik $!A$} \\
(!A \land !B)^N & \ident (!A^N \land !B^N) \\
(!A \lor !B)^N & \ident  \lnot\lnot (!A^N \lor !B^N) \\
(!A \lif !B)^N & \ident (!A^N \lif !B^N) \\
(\lforall[x][!A])^N & \ident \lforall[x][!A^N] \\
(\lexists[x][!A])^N & \ident \lnot\lnot\lexists[x][!A^N]
\end{align*}
Kolmogorov lan Glivenko
duwe versi terjemahan iki
kanggo logika proposisional;
kanggo logika predikat,
terjemahan iki ditemokake
kanthi mandiri dening
G\"odel lan Gentzen.
Kita duwe

\begin{thm}
\begin{enumerate}
\item $!A \liff !A^N$ bisa dibuktekake sacara klasik
\item Yen $!A$ bisa dibuktekake sacara klasik,
  mula $!A^N$ bisa dibuktekake sacara intuisionistik.
\end{enumerate}
\end{thm}

Saiki kita bisa mbayangake
pacelathon iki. Ahli matematika
klasik: "Aku wis mbuktekake
$!A$!" Ahli matematika
intuisionistik: "Buktimu
ora sah. Kang sejatine
wis kokbuktèkaké yaiku
$!A^N$." Ahli matematika
klasik: "Ya ora apa-apa!"
Miturut ahli matematika
klasik, ahli intuisionistik
mung mbedakake bab kang
sepele amarga loro-lorone
padha ekuivalen. Nanging
ahli intuisionistik
tetep negesake yen ana
bedane.

Gatekna yen terjemahan
ing ndhuwur mung ngenani
logika murni; terjemahan
iki ora njawab pitakonan
aksioma \emph{nonlogis}
endi kang trep kanggo
matematika klasik lan
intuisionistik, utawa
apa sesambungane aksioma
iku. Nanging teorema
ing ndhuwur bisa
dijembarake sethithik
lan menehi katrangan
kang migunani:

\begin{thm}
Yen $\Gamma$ mbuktekake $!A$ sacara klasik,
$\Gamma^N$ mbuktekake $!A^N$
sacara intuisionistik.
\end{thm}

Kanthi tembung liya,
yen $!A$ bisa dibuktekake
saka sawetara hipotesis
sacara klasik, mula
$!A^N$ bisa dibuktekake
saka terjemahan negasi
pindho hipotesis kasebut.

Kanggo nuduhake yen sawijining
ukara utawa !!{formula}
proposisional sah sacara
intuisionistik, cukup
wenehana bukti. Nanging
kepiye carane nuduhake
yen ora sah? Kanggo
kaperluan iku, kita butuh
semantik kang sah, lan
yen bisa uga jangkep.
Semantik saka Kripke
cocog kanggo kaperluan iki.

Kita bisa nglakoni kaya
ing logika klasik:
nemtokake semantik, banjur
mbuktekake kasahan lan
kelengkapan. Nanging ana
bedane kang prayoga
digatekake. Ing logika
klasik, semantike ing
sawijining pangertèn
"cetha kanthi dhewe",
wis kinandhut ing teges
panyambung-panyambunge.
Semantik Kripke bisa
diwenehi alesan intuisi,
nanging ora krasa kaya
mesthi kudu mangkono.
Kajaba iku, gagasan
!!{structure} klasik
iku lumrah ing matematika.
Mula !!a{structure}
bisa dianggep piranti
kanggo nliti logika
orde kapisan klasik,
utawa logika orde
kapisan klasik dianggep
piranti kanggo nliti
!!{structure}s matematika.
Kosok balene,
!!{structure}s Kripke
mung bisa dianggep
pambangunan logis;
ora katon nduweni
kawigaten matematika
mandiri.

!!{structure} Kripke
~$\mModel M = \tuple{W, R, V}$
kanggo basa proposisional
dumadi saka himpunan~$W$,
urutan parsial~$R$
ing~$W$ kang nduweni
!!{element} paling cilik,
lan pangaji "monoton"
variabel proposisional
marang !!{element}s
saka~$W$. Intuisine,
!!{element}s saka $W$
nggambarake "donya"
utawa "kahanan
kawruh"; unsur
$v \geq u$ nggambarake
"kahanan mbesuk kang
mungkin" saka~$u$;
dene variabel
proposisional kang
dipasangake marang~$u$
iku proposisi kang
diweruhi bener
ing kahanan~$u$.
Relasi pemaksaan
$\mSat{M}{!A}[w]$
banjur njembarake
sesambungan iki
marang sembarang
!!{formula}s ing
basa kasebut;
$\mSat{M}{!A}[w]$
diwaca "$!A$ bener
ing kahanan~$w$."
Sesambungan iki
ditemtokake kanthi
induksi kaya mangkene. Monoton tegese variabel proposisional kang wis dipasangake marang sawijining kahanan tetep dipasangake marang kabeh kahanan sabanjure; kawruh atomik ora ilang nalika kahanan maju:
\begin{enumerate}
\item $\mSat{M}{\Obj p_i}[w]$ yen lan mung yen
  $\Obj p_i$ kalebu variabel proposisional
  kang dipasangake marang~$w$.
\item $\mSat/{M}{\lfalse}[w]$.
\item $\mSat{M}{(!A \land !B)}[w]$ yen lan mung yen
  $\mSat{M}{!A}[w]$ lan $\mSat{M}{!B}[w]$.
\item $\mSat{M}{(!A \lor !B)}[w]$ yen lan mung yen
  $\mSat{M}{!A}[w]$ utawa $\mSat{M}{!B}[w]$.
\item $\mSat{M}{(!A \lif !B)}[w]$ yen lan mung yen,
  saben $w' \geq w$ lan
  $\mSat{M}{!A}[w']$ njalari
  $\mSat{M}{!B}[w']$.
\end{enumerate}
Minangka gladhen, coba
tuduhna yen $\lnot (p \land q) \lif
(\lnot p \lor \lnot q)$
ora sah sacara intuisionistik
kanthi nggawe
!!{structure} Kripke
kang dadi tuladha
kosok baline.

\end{document}
